% Midterm report skeleton -- quantum-dd-optimizer.
% STATUS: SKELETON. Contains no results. Every numerical statement must come from
% report/generated/*.tex (written by scripts from immutable results/<run_id>/ artifacts)
% and must have a row in report/evidence_map.csv (scripts/make_evidence_map.py).
% Build: scripts/build_report.sh   (latexmk -pdf -halt-on-error)
\documentclass[10pt,twocolumn]{article}

\usepackage[utf8]{inputenc}
\usepackage[margin=0.7in,columnsep=0.25in]{geometry}
\usepackage{amsmath,amssymb}
\usepackage{graphicx}
\usepackage{booktabs}
\usepackage[numbers,sort&compress]{natbib}
\usepackage[hidelinks]{hyperref}

% --- status markers (use these; never write a number by hand) -----------------
\newcommand{\notrun}{\textbf{NOT RUN}}
% run identifiers: typewriter, breakable after underscores
\newcommand{\runid}[1]{{\ttfamily\begingroup\def\_{\char`\_\allowbreak}#1\endgroup}}
\newcommand{\pending}[1]{\textbf{PENDING}\,(#1)}

% --- auto-generated numbers / status (stubs until produced by scripts) --------
% GENERATED by python scripts/make_report_artifacts.py at commit 0efe3b0; do not edit.
\newcommand{\GenIsStub}{false}
\newcommand{\GenBuildNote}{Generated from recorded runs at commit 0efe3b0}
\newcommand{\GenTestSummary}{\textbf{NOT RECORDED}}
\newcommand{\GenRunIDs}{\texttt{g3\_validation\_20261026T231608Z\_cbfe38b59354}, \texttt{benchmark\_scenario\_a\_20261026T232517Z\_1bd28ade742a}, \texttt{benchmark\_scenario\_b\_20261026T234545Z\_ad7b16de1c23}, \texttt{spectroscopy\_scenario\_a\_20261026T235246Z\_8b272601c344}}
\newcommand{\GenCommitSHA}{0efe3b0}

% GENERATED by python scripts/make_report_artifacts.py at commit de575ee; do not edit.
\newcommand{\BenchAEightBothRatioHi}{0.992}
\newcommand{\BenchAEightBothRatioLo}{0.980}
\newcommand{\BenchAEightPhaseRatioHi}{1.000}
\newcommand{\BenchAEightPhaseRatioLo}{0.999}
\newcommand{\BenchAEightSel}{XY8}
\newcommand{\BenchAEightSelEps}{1.01\times10^{-2}}
\newcommand{\BenchAEightTimingRatioHi}{0.203}
\newcommand{\BenchAEightTimingRatioLo}{0.203}
\newcommand{\BenchAFourBothRatioHi}{1.011}
\newcommand{\BenchAFourBothRatioLo}{0.965}
\newcommand{\BenchAFourPhaseRatioHi}{1.002}
\newcommand{\BenchAFourPhaseRatioLo}{1.001}
\newcommand{\BenchAFourSel}{XY4}
\newcommand{\BenchAFourSelEps}{1.06\times10^{-2}}
\newcommand{\BenchAFourTimingRatioHi}{0.507}
\newcommand{\BenchAFourTimingRatioLo}{0.500}
\newcommand{\BenchASixteenBothRatioHi}{0.994}
\newcommand{\BenchASixteenBothRatioLo}{0.982}
\newcommand{\BenchASixteenPhaseRatioHi}{0.998}
\newcommand{\BenchASixteenPhaseRatioLo}{0.996}
\newcommand{\BenchASixteenSel}{XY8}
\newcommand{\BenchASixteenSelEps}{1.00\times10^{-2}}
\newcommand{\BenchASixteenTimingRatioHi}{0.071}
\newcommand{\BenchASixteenTimingRatioLo}{0.071}
\newcommand{\FloorA}{9.90\times10^{-3}}
\newcommand{\GthreeMaxZ}{0.95}
\newcommand{\GthreeNtraj}{64000}


\title{Model-Based Dynamical Decoupling under Hardware Timing Constraints\\for Idle Superconducting Qubits\\[0.5ex]
  \large Midterm report --- synthetic-simulation study (no hardware data)}
\author{TEAM AUTHOR LIST --- PENDING CONFIRMATION\\
  \small Columbia University (pending confirmation)}
\date{\GenBuildNote}

\begin{document}
\maketitle

\begin{abstract}
% Abstract. Numbers are generated macros (report/generated/numbers.tex).
An idle superconducting qubit waiting during a neighbour's mid-circuit measurement loses
coherence to stochastic dephasing and to deterministic frequency shifts. We ask whether
model-based optimization of dynamical-decoupling (DD) pulse timings and phases with CMA-ES
outperforms strong canonical sequences. To answer it fairly, we pre-registered a synthetic
transmon-like model, seeds, baselines and a training-only baseline-selection rule, and scored
every arm on independent held-out noise. A finite-pulse Monte-Carlo simulator agrees with
filter-function theory in five predetermined configurations (largest $|z|=\GthreeMaxZ$).
Synthetic Ramsey/Hahn/CPMG spectroscopy recovers a predictive noise model (held-out RMSE
$\SpecGaussFullRmseRatio\times$ the shot-noise floor), but the spectrum is not uniquely
identified. In the pre-registered Scenario~A, no CMA-ES mode beat the selected XY4/XY8 baseline,
whose infidelity ($\BenchASixteenSelEps$ at $n=16$) is already at the $T_1$-plus-white-noise
floor ($\FloorA$). With a deterministic Stark-like shift and virtual-$Z$ correction, the only
consistent gain was small (ratio $\BenchBFourBothRatioLo$--$\BenchBFourBothRatioHi$ at $n=4$).
Optimizing on fitted models, or with more headroom, did not help, and training gains did not
generalize. All results are synthetic; no quantum hardware was used.

\end{abstract}

\section{Introduction and project bounds}
\label{sec:intro}
% ~0.75 page. Motivation (IQCC/Columbia), idle qubit during mid-circuit measurement (MCM),
% metrics, hypotheses and scope.

\paragraph{Motivation} An idle spectator qubit must preserve coherence while
other qubits are measured and classical feedback is processed. Dynamical
decoupling can suppress suitable low-frequency dephasing, but finite pulse
errors, energy relaxation, and broadband fluctuations may limit improvement.
We investigate whether model-specific pulse optimization provides gains
beyond established XY4/XY8 sequences when the available error budget is
already near its physical floor. The present analysis is a synthetic
computational study; superconducting-hardware validation is a separate
planned experiment.

\paragraph{Problem} An idle qubit accumulates phase errors from stochastic dephasing and, in
dynamic circuits, from deterministic frequency shifts during neighbouring readout. We ask
whether model-based optimization of DD sequences can reduce idle-window infidelity relative to
strong canonical baselines under a fixed pulse budget and identical constraints.

\paragraph{Research questions}
\begin{description}
  \item[RQ1] Agreement of Monte Carlo coherence with analytic filter-function predictions under
    matched PSD conventions.
  \item[RQ2] Predictive (held-out) accuracy of noise parameters fitted from synthetic
    Ramsey, Hahn and CPMG data.
  \item[RQ3] Whether joint optimization of pulse timings and phases beats strong DD baselines in
    Scenario~A (stochastic dephasing) on held-out noise.
  \item[RQ4] Whether any gain over calibrated virtual-$Z$ and standard DD survives uncertainty
    in readout-induced detuning (Scenario~B).
  \item[RQ5] Robustness to pulse imperfections and model mismatch. Additional
    drift, Stark rise-time and control-quantization tests remain future work.
\end{description}

\paragraph{Scope} Synthetic simulations only. Scenarios, seeds, baselines, the
baseline-selection rule and the metrics were pre-registered in the repository before any
optimization run; deviations are logged. A zero gain or a regression counts as a valid outcome. Metric: average gate infidelity against the corrected target
(Sec.~\ref{sec:model}).

\section{Related work}
\label{sec:related}
% ~0.75 page. Each sentence maps to a LIT-xx claim in report/claims.json.
% Primary arXiv identities were checked in a later audit; consult docs/CITATION_AUDIT.md.

% LIT-01
Dynamical decoupling of open quantum systems was introduced by \citet{Viola1999}.
% LIT-02
\citet{Uhrig2007} proposed optimized $\pi$-pulse sequences for keeping a qubit coherent,
% LIT-03
and \citet{Cywinski2008} analysed how to enhance dephasing times in superconducting qubits.
% LIT-04
\citet{Biercuk2009} studied optimized DD in a model quantum memory.
% LIT-05, LIT-06
Biercuk, Doherty, and Uys described DD optimization as a filter-design
problem \citep{BiercukDohertyUys2011}, while Green and coauthors developed
higher-order filter methods for driven gates \citep{GreenUysBiercuk2012}.
These approaches directly motivate deterministic model-based objectives,
which differ from the finite-trajectory CMA-ES objective used here.
DD has also been used for noise spectroscopy: \citet{Bylander2011} with a superconducting flux
qubit and \citet{Alvarez2011} for colored noise.
% LIT-07
Robust DD was studied by \citet{Souza2011},
% LIT-08
and \citet{Ezzell2023} surveyed DD performance on superconducting qubits.
% LIT-09, LIT-10, LIT-11
Empirical or hardware-tailored learning approaches include \citet{Rahman2024}, \citet{Tong2025}
and \citet{Tong2026}.

\paragraph{Empirical versus model-based}
% LIT-10 (title only)
\citet{Rahman2024}, \citet{Tong2025}, and \citet{Tong2026} study
empirical or hardware-tailored dynamical decoupling. In particular,
\citet{Rahman2024} optimizes rotational gates for superconducting
hardware, while \citet{Tong2026} studies empirical suppression in dynamic
circuits. The present approach instead uses a fitted or assumed synthetic
noise model as the optimization objective. Such a model-based approach
reduces direct hardware calls but depends on model fidelity.

\paragraph{Positioning}
% No LIT claim: describes this project only. Novelty claim deliberately limited.
This project is a model-based, synthetic-only study: sequences are optimized and evaluated
against a simulated noise model, and no quantum hardware has been used.
Its contribution is limited to a pre-registered, held-out synthetic benchmark of model-based DD
optimization with explicit noise floors.
Its main finding is negative: under the pre-registered synthetic model, CMA-ES did not
provide a practically meaningful generalizable improvement over XY4/XY8 on held-out noise (Sec.~\ref{sec:results}).
A synthetic benchmark of this kind does not test whether its conclusions transfer to hardware.

\section{Physical model and conventions}
\label{sec:model}
% ~0.75 page. Conventions follow docs/PHYSICS_CONVENTIONS.md (integration-lead owned).
% Definitions only. Derivations: docs/PHYSICS.md.

We use $\hbar=1$, time in s, angular frequency and detuning in rad/s and angles in rad. In the
rotating frame of the drive,
\begin{multline}
  H(t) = \tfrac{1}{2}\bigl[\delta(t) + d(t) + \delta_{\mathrm{off}}\bigr]\sigma_z\\
       + \tfrac{1}{2}\Omega(t)\bigl[\cos\varphi(t)\,\sigma_x + \sin\varphi(t)\,\sigma_y\bigr],
\end{multline}
where $\delta(t)$ is zero-mean stationary noise, assumed \emph{Gaussian} unless stated
otherwise, $d(t)$ a known deterministic detuning profile (e.g.\ during neighbouring readout),
and $\delta_{\mathrm{off}}$ a static calibration offset.

\paragraph{Spectrum and filter} With autocorrelation
$C(\tau)=\langle\delta(t+\tau)\delta(t)\rangle$ we use the \emph{two-sided} PSD
$S_2(\omega)=\int_{-\infty}^{\infty} C(\tau)e^{-i\omega\tau}\,d\tau$ in rad$^2$/s.
For the switching function $y(t)$ (toggling-frame $z$ component, $\pm1$ between ideal $\pi$
pulses),
\begin{gather}
  Y(\omega)=\int_0^T y(t)e^{i\omega t}\,dt,\\
  \chi=\operatorname{Var}\phi=\frac{1}{2\pi}\int_{-\infty}^{\infty} S_2(\omega)\,|Y(\omega)|^2\,d\omega,
\end{gather}
and, under the Gaussian assumption, $W=e^{-\chi/2}$ with pure-dephasing average gate infidelity
$r=(1-W)/3$. Low- and high-frequency cutoffs of $S_2$ are stated explicitly for every model.

\paragraph{Energy relaxation} $T_1$ is modelled as amplitude damping with
$\gamma=1-e^{-t/T_1}$ per free segment; channels are composed as Pauli transfer matrices, never
by adding infidelities. Amplitude damping alone has average infidelity
$r_{\rm AD}=(2+\gamma-2\sqrt{1-\gamma})/6\approx\gamma/3$ (derived in \repopath{docs/PHYSICS.md} and
checked numerically against Haar and six-state averages).

\paragraph{Target and frame correction} Fidelity is evaluated against
$U_{\mathrm{target}}=R_z(\theta_{\mathrm{vz}})U_0$, where $U_0$ is the ideal net control
unitary of the sequence and $\theta_{\mathrm{vz}}$ is a virtual-$Z$ correction of the known
$d(t)$ (zero unless frame correction is allowed). Average gate fidelity of a channel with PTM
$R$ is $F_{\mathrm{avg}}=(\operatorname{Tr}R+2)/6$ evaluated on $R_{\mathrm{target}}^{\mathsf T}R$.
For finite pulses we use the generalized (three-component toggling-frame) filter function, so
$W=e^{-\chi/2}$ is exact only for instantaneous pulses and leading-order otherwise; the
Monte-Carlo simulator needs neither approximation.
Finite pulse widths, minimum spacing and timing-grid constraints are enforced for all sequences.

\section{Methods (planned; implementation status pending)}
\label{sec:methods}
% ~1.0 page. Describe ONLY what the code actually implements once it exists; until then each
% item is labelled with its status.

\paragraph{Synthetic device and noise.} Gaussian PSD models (band-limited $1/f^\alpha$,
Lorentzian, white) and a non-Gaussian telegraph process for misspecification tests.
Status: \pending{ddopt.noise, Agent A}.

\paragraph{Finite-pulse simulator.} Noise-averaged channel built from propagation through finite
pulses and free intervals, graded against the corrected target. Status:
\pending{ddopt.simulator, Agent B}.

\paragraph{Spectroscopy and held-out fits.} Synthetic Ramsey/Hahn/CPMG data with shot noise and
SPAM; fitted parameters are evaluated on disjoint held-out pulse counts and idle times.
Status: \pending{ddopt.spectroscopy, Agent C}.

\paragraph{Optimization and baselines.} CMA-ES \citep{Hansen2016} (LIT-12) over pulse timings
and/or phases; comparison with free evolution, Hahn, CPMG, XY4, XY8, UDD and (Scenario~B)
virtual-$Z$ at equal $n$, $T$, pulse width and correction privileges. The baseline is selected
on training draws by a pre-registered rule; final scores use independent test seeds.
Status: \pending{ddopt.optimizer / ddopt.evaluation, Agent D}.

\paragraph{Reproducibility.} Each run writes an immutable directory
\texttt{results/<run\_id>/} with a manifest (UTC time, git SHA and dirty flag, platform, Python
and dependency versions, config hash, seeds, command) and SHA-256 checksums
(\texttt{ddopt.provenance}). Tables and numbers in this report are generated from those files;
the claim-to-artifact ledger is \texttt{report/evidence\_map.csv}. Test status:
\GenTestSummary.

\paragraph{Hardware plan.} A QUA adapter is prepared as a template only; no hardware execution
is authorized or claimed. Status: \pending{docs/HARDWARE\_PROTOCOL.md, Agent E}.

\section{Experiments and results}
\label{sec:results}
% ~1.25--1.75 pages once results exist. Only executed experiments may be reported, each with
% run ID, seeds and a row in report/evidence_map.csv. Figures must be produced by scripts from
% stored results; no hand-made or placeholder plots.

\textbf{NOT RUN.} No experiment has been executed for this draft (claim PROV-02). The tables
below are stubs from \texttt{report/generated/} and contain no values. Recorded run IDs:
\GenRunIDs.

% STUB -- NOT GENERATED FROM ANY RUN. Will be overwritten by a script that reads the
% sequence-feasibility results (n, T, pulse width, guard spacing, timing-grid rounding).
\begin{table}[t]
  \centering
  \caption{Timing feasibility of the sequence configurations. \textbf{NOT RUN}: stub table,
    no configurations evaluated.}
  \label{tab:feasibility}
  \begin{tabular}{@{}llllll@{}}
    \toprule
    $n$ & $T$ & Pulse width & Guard & Grid rounding & Feasible \\
    \midrule
    --- & --- & --- & --- & --- & \textbf{NOT RUN} \\
    \bottomrule
  \end{tabular}
\end{table}

% STUB -- NOT GENERATED FROM ANY RUN. Will be overwritten by a script that reads
% results/<run_id>/payload.json (equal-n, equal-T benchmark with confidence intervals).
\begin{table}[t]
  \centering
  \caption{Equal-$n$, equal-$T$ benchmark on held-out test seeds. \textbf{NOT RUN}: stub table,
    no values computed.}
  \label{tab:benchmark}
  \begin{tabular}{@{}llll@{}}
    \toprule
    Sequence & $n$ & Test infidelity (95\% CI) & Run ID \\
    \midrule
    Free evolution    & --- & \textbf{NOT RUN} & --- \\
    Canonical DD (Hahn, CPMG, XY4, XY8, UDD) & --- & \textbf{NOT RUN} & --- \\
    Virtual-$Z$ (Scenario B) & --- & \textbf{NOT RUN} & --- \\
    CMA-ES optimized  & --- & \textbf{NOT RUN} & --- \\
    \bottomrule
  \end{tabular}
\end{table}


\paragraph{Planned exhibits (all \notrun).}
\begin{enumerate}
  \item True versus fitted PSD with units and credible interval (RQ2).
  \item Held-out Ramsey/Hahn/CPMG decay curves with uncertainty (RQ2).
  \item Monte Carlo versus filter-function agreement across $\ge 3$ predetermined configurations
    (RQ1, gate G3).
  \item Equal-$n$, equal-$T$ benchmark with confidence intervals (Table~\ref{tab:benchmark}; RQ3).
  \item Optimization training versus held-out convergence (RQ3).
  \item Robustness to readout timing/amplitude, PSD misspecification and pulse imperfections
    (RQ4, RQ5), including negative findings.
\end{enumerate}
Improvement factors, if any, will be reported as
$\epsilon_{\mathrm{baseline}}/\epsilon_{\mathrm{optimized}}$ together with both raw infidelities.

\section{Discussion, limitations and roadmap}
\label{sec:discussion}
% ~0.75 page.

\paragraph{Discussion.} \pending{depends on results; none available}.

\paragraph{Limitations (known in advance).}
\begin{itemize}
  \item All planned evidence is synthetic; a synthetic model can be optimized against itself, so
    gains may not transfer to hardware (model mismatch).
  \item The filter-function infidelity relies on a Gaussian stationary noise assumption;
    non-Gaussian (e.g.\ telegraph) noise is only approximated.
  \item Spectroscopy is limited by finite shots and SPAM; fitted spectra are not uniquely
    identifiable.
  \item No quantum-hardware execution has been performed (claim PROV-01); access to IQCC
    hardware is unconfirmed, and SDK compilation, if achieved, is not hardware validation.
\end{itemize}

\paragraph{Roadmap.} \pending{to be written from the actual state of the implementation}.

\section*{Author Contributions and Disclosure}
The authors are Christian De Zabala F., Madeline Feldner, Zarar Haider, and
Udaya Vijay Anand. Specific individual contributions are not assigned here
because the four authors have not yet confirmed those attributions.

\textbf{Use of AI tools:} AI coding assistance was used in repository scaffolding,
implementation, and manuscript preparation. The scientific interpretations and
attribution require author review. All reported measurements are synthetic,
and no IQCC hardware result is claimed.


\bibliographystyle{plainnat}
\bibliography{refs}

\end{document}
